\chapter{Stage-Gate and Development Strategy}

\section{Session Overview}

This session asks your group to decide how the project should be managed before detailed design work begins. In practice, a product team does not only choose \emph{what} to develop; it also decides \emph{how} the work will be organised, reviewed, and coordinated. These choices affect risk, speed, communication, and the quality of later design decisions.

By the end of this session, your group should be able to:
\begin{enumerate}
    \item identify a suitable product development process for your hydroponics project;
    \item select a development process flow that matches the project's complexity and uncertainty;
    \item distinguish between departmental structure and organisational structure; and
    \item justify a recommended team structure for the project.
\end{enumerate}

\section{Why a Stage-Gate Approach Matters}

A \textbf{stage-gate process} breaks product development into major stages separated by review points, or \emph{gates}. At each gate, the team checks whether the project should continue, be revised, or stop. This approach is useful because it forces teams to review evidence before moving forward.

For APID, the stage-gate mindset helps your group do three things:
\begin{itemize}
    \item confirm that the chosen opportunity still fits user needs and project goals;
    \item identify technical or market risks early; and
    \item coordinate decisions across engineering, design, communication, and planning tasks.
\end{itemize}

\section{Key Terms}

This section separates terms that are often confused.

\begin{itemize}
    \item \textbf{Product development process}: the overall type of development strategy, such as market-pull or technology-push.
    \item \textbf{Development process flow}: the sequence pattern used to execute the work, such as a generic linear flow or a spiral cycle.
    \item \textbf{Departmental structure}: the functions that must contribute to the work, such as marketing, engineering, manufacturing, or quality.
    \item \textbf{Organisational structure}: how authority and responsibility are arranged, such as a functional team, project team, lightweight structure, or heavyweight structure.
\end{itemize}

\section{Development Process Types}

Table~\ref{tab:session2-process-types} summarises common process types and shows how they may relate to hydroponic products.

\begin{table}[htbp]
\centering
\caption{Summary of common product development process types}
\label{tab:session2-process-types}
\renewcommand{\arraystretch}{1.2}
\begin{tabularx}{\textwidth}{|>{\raggedright\arraybackslash}p{0.2\textwidth}|>{\raggedright\arraybackslash}X|>{\raggedright\arraybackslash}p{0.25\textwidth}|}
\hline
\textbf{Process type} & \textbf{When it fits} & \textbf{Hydroponics-related example} \\
\hline
Generic or market-pull & Starts from a clear user need and then selects technology to meet it. & A grower-facing cleaning tool developed to solve a maintenance problem reported by users. \\
\hline
Technology-push & Starts from a new technology and then looks for a valuable application. & A new low-cost sensor package adapted for nutrient monitoring. \\
\hline
Platform & Builds a new product around an existing subsystem or architecture. & A modular rack system designed around an existing pump-and-channel platform. \\
\hline
Process-intensive & Product characteristics are strongly shaped by the production or operating process. & A nutrient-mixing product whose performance depends on controlled dosing and mixing behaviour. \\
\hline
Customized & Creates tailored variations from a known product family. & A hydroponic support frame adapted to different farm or balcony layouts. \\
\hline
High-risk & Involves significant market or technical uncertainty and requires early testing. & A novel automated dosing system with uncertain reliability in real growing conditions. \\
\hline
Quick-build & Uses repeated build-test cycles to refine the concept rapidly. & A low-cost prototype enclosure or fixture refined through repeated physical iteration. \\
\hline
Product-service system & Combines a physical product with a supporting service or workflow. & A monitoring kit bundled with maintenance guidance and data reporting. \\
\hline
Complex system & Requires several interacting subsystems to work together reliably. & A controlled hydroponic unit integrating sensing, pumping, structure, and user interface. \\
\hline
\end{tabularx}
\end{table}

When your group chooses a process type, focus on the product's dominant challenge. For example, if the main uncertainty is a novel sensor technology, a technology-push or high-risk framing may fit better than a generic market-pull process.

\section{Development Process Flows}

After selecting the broad process type, choose the flow that best suits the project execution.

\begin{itemize}
    \item \textbf{Generic flow}: a mostly sequential path from planning to concept development, design, testing, and refinement. This works well for low-to-moderate uncertainty.
    \item \textbf{Spiral flow}: repeated cycles of planning, prototyping, review, and revision. This is useful when technical learning must happen through iteration.
    \item \textbf{Complex-system flow}: parallel work across multiple subsystems followed by integration. This is appropriate when the product contains many interacting modules.
\end{itemize}

If your product includes electronics, fluid handling, and structure that must work together, a spiral or complex-system flow is often more realistic than a single-pass linear flow.

\section{Departmental Structure}

This section focuses on \emph{who} needs to contribute. A product development team commonly draws from:
\begin{itemize}
    \item marketing or user research;
    \item engineering and technical analysis;
    \item manufacturing or fabrication planning;
    \item quality and testing;
    \item purchasing or supply coordination;
    \item finance or cost planning; and
    \item project management.
\end{itemize}

For APID, your group does not need to create a large company chart. Instead, decide which functions are most important for your project and explain why. A sensor-heavy hydroponic product may need stronger emphasis on instrumentation and testing, whereas a structural product may need more attention to materials and fabrication.

\section{Organisational Structure}

This section focuses on \emph{how} authority is arranged.

\begin{table}[htbp]
\centering
\caption{Common organisational structures for product development teams}
\label{tab:session2-organisational-structures}
\renewcommand{\arraystretch}{1.2}
\begin{tabularx}{\textwidth}{|>{\raggedright\arraybackslash}p{0.18\textwidth}|>{\raggedright\arraybackslash}X|>{\raggedright\arraybackslash}p{0.26\textwidth}|}
\hline
\textbf{Structure} & \textbf{Main characteristic} & \textbf{Typical use in APID terms} \\
\hline
Functional & Work stays mainly within specialised roles or disciplines. & Suitable when the project is simple and roles are stable. \\
\hline
Project & The team is organised mainly around the project itself. & Suitable when the group needs fast decisions and strong coordination. \\
\hline
Lightweight & A project lead coordinates, but most authority remains with functional roles. & Suitable when the work is shared but not highly integrated. \\
\hline
Heavyweight & The project lead has stronger authority across functions and priorities. & Suitable when integration, speed, and cross-functional decisions matter greatly. \\
\hline
\end{tabularx}
\end{table}

Most APID groups will choose either a project structure or a lightweight structure because both support shared responsibility without requiring a full corporate hierarchy.

\section{Pre-Class Preparation}

Before class, each group member should prepare short notes covering the following decisions:
\begin{enumerate}
    \item the most suitable development process type for the selected hydroponic product;
    \item the most suitable development process flow;
    \item the most important departments or functions needed for the project; and
    \item the organisational structure that would best coordinate the work.
\end{enumerate}

For each decision, write one short justification based on project complexity, uncertainty, user needs, technical risks, or team coordination.

\section{In-Class Activity: Consolidate the Development Strategy}

In class, compare the individual proposals and agree on one team recommendation for each category.

\subsection{Suggested Decision Steps}
\begin{enumerate}
    \item compare each member's proposed process type and identify the strongest justification;
    \item agree on the development process flow that best matches the level of iteration or subsystem integration required;
    \item decide which departments or functions are essential to the project;
    \item choose the organisational structure that supports those functions most effectively; and
    \item document the final decision in a short report for ITEL submission if requested.
\end{enumerate}

\subsection{Decision Matrix Template}

Use Table~\ref{tab:session2-decision-matrix} if your group wants a structured way to compare options.

\begin{table}[htbp]
\centering
\caption{Simple decision matrix for Session 2 selections}
\label{tab:session2-decision-matrix}
\renewcommand{\arraystretch}{1.2}
\begin{tabularx}{\textwidth}{|>{\raggedright\arraybackslash}p{0.22\textwidth}|>{\raggedright\arraybackslash}p{0.18\textwidth}|>{\raggedright\arraybackslash}p{0.16\textwidth}|X|}
\hline
\textbf{Decision area} & \textbf{Selected option} & \textbf{Confidence (1--5)} & \textbf{Justification} \\
\hline
Development process type &  &  &  \\
\hline
Development process flow &  &  &  \\
\hline
Departmental structure &  &  &  \\
\hline
Organisational structure &  &  &  \\
\hline
\end{tabularx}
\end{table}

\section{Summary and Next Steps}

By the end of this session, your group should have one agreed development strategy covering process type, process flow, departmental structure, and organisational structure. Keep this summary because it will influence later decisions about planning, customer needs, concept development, and project communication.
